\chapter{Conclusion}
\label{ch:conclusion}
This thesis now distinguishes a primary A-CDM empirical analysis from an earlier auxiliary simulation study. The primary collection contains 1,322 departure records from five saved dates. Source checks establish CSV/JSON consistency and traceability. Of 1,286 eligible actual takeoffs, 1,257 have a unique OpenSky callsign/time candidate. The close timestamp correspondence supports the departure comparison for that subset, subject to proxy and timezone limitations.

For RQ1, the A-CDM observations support departure-level descriptions and checks on selected attributes. They do not provide arrival/block milestones or internal service observations needed for full activity calibration. OpenSky validation cannot supply those missing fields while also claiming that all inputs originate in A-CDM.

For RQ2 and RQ3, the auxiliary OCPN/Promenade work demonstrates analysis of shared-resource waiting and selected capacity scenarios. Its 60 replications and numerical findings remain unchanged in the appendices. The 10,940-row input cohort originates in OpenSky, and its generated outputs remain conditional model evidence. These results do not establish an A-CDM-only calibrated BER turnaround model or optimal real staffing.

\section{Future Work}
The next step is to agree on the empirical scope with the supervisor and obtain reliable arrival/block milestones if full turnaround calibration remains the objective. Further work must also correct local-time movement policy, separate long stays and evaluate matching observable quantities on independently selected dates. Terminal/remote comparisons and controlled inter-flight delay propagation remain necessary for a fuller answer to RQ2.

The completed work provides an auditable A-CDM primary dataset, a reproducible departure cross-check and a documented exploratory simulation. Its value depends on preserving these distinctions rather than increasing the apparent evidence by relabeling reconstructed or simulated records.
